\section{Corrections and status revisions for the manuscript}
\label{editorial:sec}
The source's discovery chapter already gives a careful distinction
between experimental evidence, exact relation matrices, and numerical
period identities. A few earlier passages should be brought into line
with that standard. The proposed patch is restricted to the pinned
source and is supplied for review; it has not been applied to the
remote repository.

\subsection{A definite bibliographic correction}
The plastic-base discussion in
\path{chapters/03-algebraic.tex} describes a sequel of
Abouzahra--Lewin as never published. The sequel was published:
M.~Abouzahra, L.~Lewin, and Xiao Hongnian,
\emph{Polylogarithms in the field of omega (a root of a given cubic):
Functional equations and ladders},
\emph{Aequationes Mathematicae} \textbf{33} (1987), 23--45
\cite{AbouzahraLewinXiao1987}.
The publisher records the issue date as February 1987 and the DOI as
\texttt{10.1007/BF01836149}.
Its abstract distinguishes analytically derived lower-order results
from numerically found higher-order ladders. Restoring this reference
therefore improves both attribution and proof-status tracking.

The later literature also includes Gangl's functional equations and
proved ladders in weights six and seven \cite{Gangl2013}.
Those statements use a single-valued modified polylogarithm; the
projection is the real part in odd weight and the imaginary part in
even weight. In particular, an even-weight modified identity at real
arguments cannot simply be read as an identity for ordinary real
\(\Li_{2k}\). The correction patch makes this normalization distinction
explicit.

\subsection{Numerical confirmation is a distinct status}
The golden-ladder discussion uses headings and phrases such as
\emph{conjectures, verified}, \emph{also hold}, and \emph{are correct}
beside finite-precision residual checks. A small computed residual
establishes numerical consistency with a formula. By itself it does
not establish an identity of periods. The proposed edits say
\emph{numerically confirmed} where that is the evidence actually
provided, and retain analytic or exact status when a proof is present.

This is a local editorial correction, not a claim that every displayed
historical ladder lacks a proof in the literature. A formula with a
published proof should be matched to that proof, including its
normalization and branches. A proposed specialization of a functional
equation should include the exact substitution and residual
cancellation. In particular, citing a family of Kummer equations
without supplying or identifying the required specialization is less
specific than a replayable derivation.

The new seed theorem sharpens the source's support question but does
not remove this distinction. It proves that no new positive even
golden power adds a multiplicative seed. It does not prove the
constants in each higher ladder or a maximal ladder weight of nine.
The remaining question concerns compatibility with functional equations,
products, and regulator data on an already complete seed space.

\subsection{Restoring a missing canonical ladder definition}
The canonical elimination chain in the same source uses \(L_{12}\)
without defining it. Its logarithmic tail also contains factorials
with negative indices at the low orders where vanishing is discussed.
The proposed repair omits those summands and reconstructs the missing
notation from the first displayed weight-five formula.
Put \(\ell=\log\rho\) and, for \(n\ge1\), define
\begin{align}
 \mathrm{tw}_{12}(n)={}&-\frac{13}{48}\frac{\ell^n}{n!}
 +\frac1{48}\zeta(2)\frac{\ell^{n-2}}{(n-2)!}
 -\frac{19}{1728}\zeta(4)\frac{\ell^{n-4}}{(n-4)!},
 \label{editorial:eq:tail}\\
 L_{12}(n)={}&\frac{\Li_n(\rho^{12})}{12^{n-1}}
 -\frac32\frac{\Li_n(\rho^6)}{6^{n-1}}
 -\frac{\Li_n(\rho^4)}{4^{n-1}}
 +\frac{11}{48}\frac{\Li_n(\rho^2)}{2^{n-1}}
 +\mathrm{tw}_{12}(n).
 \label{editorial:eq:L12}
\end{align}
The negative-factorial convention applies in \eqref{editorial:eq:tail}.
If \(E_1,E_2,E_3\) are the left-minus-right residuals of the
three displayed source weight-five formulas, the exact rational
normalizations are
\begin{align}
 L_{12}(5)-\frac{67}{6912}\zeta(5)
    &=\frac{E_1}{15\cdot12^4},\notag\\
 L_{20}(5)-\frac{201}{10000}\zeta(5)
    &=\frac{E_2}{9\cdot20^4},\label{editorial:eq:normalization}\\
 L_{24}(5)+\frac{1541}{110592}\zeta(5)
    &=\frac{E_3+(64/3)E_1}{-15\cdot24^4}.\notag
\end{align}
Thus the canonical \(L_{24}\) requires a recombination before
rescaling. These equalities follow by exact rational coefficient
comparison, and restore a reviewable meaning to the source's
elimination chain. They do not assert \(E_j=0\) without a proof
of the underlying evaluations.

The replay program checks all three transformations exactly and
compares \(\Li_n\) from \texttt{mpmath} with a separate 600-term
defining series at 160 working digits. The largest recorded canonical
residual over orders one through five is below
\(1.12\times10^{-161}\). The series has an analytic truncation
bound; floating-point rounding is not interval-enclosed.
The notation repair and the numerical evaluation therefore retain
their different logical statuses.

\subsection{Limits that should remain explicit}
The fixed-\(m\) moment theorem in the source is sound in its stated
regime. The factor \(e^{c e^{-s}}\) in
Theorem~\ref{moment:thm:transition} explains why extending that regime
without a uniform analysis would be incorrect. The source's broader
question about simultaneous large exponents is only partially resolved:
the present strip is near \(n\sim m\log m\), and does not cover every
fixed ratio \(n/m\).

For real-order depth-two functions, the equality \(a+b=1\) changes the
meaning of a unit-circle value. The coefficients then tend to \(1/a\),
so individual terms of the boundary series do not tend to zero.
The signed-measure continuation gives well-defined analytic values
away from the cut; ordinary series convergence must not be asserted.

Finally, the \(S_4\) identity is already proved in the pinned source.
Its \(S_6\) identity has total weight seven and remains conjectural.
The source also already proves that two displayed \(S_6\) coordinate
forms are equivalent. That exact coordinate change is not a proof
that either conjectural residual vanishes.

\section{Further research questions and proposed relations}
\label{future:sec}
The following program separates promising conjectures from questions
for which even the correct global statement remains uncertain.
The priority order reflects how directly each task builds on a
proved result in this article.

\subsection{From complete seeds to analytic golden ladders}
The complete even-power alphabet is now
\[
 \{\rho^2,\rho^4,\rho^6,\rho^8,\rho^{10},
       \rho^{12},\rho^{20},\rho^{24}\}.
\]
The next task is to obtain explicit analytic certificates for the
source's higher golden formulas on this finite set. A useful
certificate would list the functional equations used, the algebraic
substitutions, every lower-weight product contribution, and the
constant fixed by evaluation at a regular point. The six integral
seed vectors provide canonical input rows, and the exterior-symbol
calculation supplies an exact first obstruction at every weight.

Three questions are especially concrete. Which of the four even
symbol directions extend through each weight? Can the regulator or
coaction constraints be expressed in a finite sequence of exact
matrices? If extension ceases in a specified ladder class, can one
prove that failure inside that class? A numerical absence of another
ladder, or the completeness of its argument list, is not such a proof.

\subsection{Classifying all quadratic bases with nontrivial seeds}
The uniform cutoff reduces any positive real quadratic unit to at
most twenty-four indices. One can therefore seek a classification in
terms of its trace, norm, and power within the unit group. For a fixed
base, ideal factorization followed by an integral nullspace and unit
calculation gives a finite algorithm. The stronger family problem is
to determine all bases for which the nullspace is nonzero, without
enumerating an unbounded range of traces.

The golden ranks \(6,4,1,1,0,\ldots\) show why retaining the base, not
merely its number field, matters. Extending this classification to
cubic or quartic units will require a different effective
primitive-divisor analysis and careful treatment of the larger unit
rank. The plastic examples are a particularly relevant first case
because their published functional equations are already available
\cite{AbouzahraLewinXiao1987}.
Campbell's more recent construction of dilogarithm families with
unboundedly many terms in the defining base polynomials
\cite{Campbell2024} supplies a complementary direction in which the
degree of the base is allowed to grow. A useful comparison would
separate that variable-degree construction from the fixed-degree
primitive-divisor bounds used here.

\subsection{A proposed strict angular monotonicity}
Theorem~\ref{real:thm:arc} proves \(x'(s)>0\), but the angle also
depends on the changing radius. The stronger statement is the
following plausible continuation.

\begin{conjecture}[Radial angular monotonicity]
\label{future:conj:angle}
For every \(a,b>0\) with \(a+b\ge1\),
\[
 \frac{d}{dR}\theta_{a,b}(R)<0
 \qquad(0<R\le1).
\]
Equivalently, the analytic abscissa from
Theorem~\ref{real:thm:arc} satisfies
\[
 2s\,x_{a,b}'(s)>x_{a,b}(s)\qquad(0<s\le1).
\]
\end{conjecture}

The small-radius expansion proves the assertion for sufficiently
small \(R\), for each fixed pair \(a,b\).
It holds exactly when \(a=b=1\), and the recorded numerical samples
show decreasing angles at the four tested radii for each of four
parameter pairs. That is limited evidence, not a proof over the
parameter domain. The existing sign argument establishes positivity
of \(G_x\) and negativity of \(G_s\); it does not establish the
stronger comparison \(-2sG_s>xG_x\) required here.
A proof may need a further inequality for the cumulative signed
measure, or a counterexample may reveal the correct restricted domain.

Separate questions concern monotonicity in \(a\) and \(b\).
The strictly decreasing density-crossing location in
Corollary~\ref{real:cor:order-motion} does not settle them. A
likelihood-ratio comparison for the cumulative measures would give
more useful information than comparing the densities' zeros alone.

\subsection{Below the finite-measure threshold and at its corners}
For \(a+b<1\), the coefficients grow like \(n^{1-a-b}/(1-b)\).
Finite signed measures are therefore impossible, but the generating
function and its analytic continuation remain meaningful.
Can a controlled derivative or fractional derivative of a measure
replace the representation used here? Does the unique angular zero
survive throughout that larger domain, and what happens to it as
\(a+b\downarrow0\)? These are open questions in this continuation,
not consequences of the signed-measure theorem.

At the critical boundary, Section~\ref{real:sec:boundary} fixes \(b\) in
\((0,1)\). Joint limits \(b\downarrow0\) or \(b\uparrow1\) can change
the scale and the constants. Determining a uniform description in
those corners would connect power-law density crossings with
logarithmic regimes. The proved exponential concentration law also
suggests using its rescaled coordinate for quadrature, with explicit
error bounds that remain stable as \(a+b\downarrow1\).

\subsection{The full two-exponent moment problem}
The integral obeys the exact symmetry \(M_{n,m}=M_{m,n}\).
The present transition describes one highly asymmetric region, and
reflection gives the opposite region. A complete asymptotic theory
should connect these regions with the interior saddle for bounded
positive \(n/m\). A starting point is the phase
\[
 \Phi_\alpha(x)=\alpha\log f(x)+(1-\alpha)\log f(1-x),
 \qquad \alpha=\frac{n}{n+m}.
\]
Before asserting a uniform saddle formula, one should prove the
number and nondegeneracy of its maxima as \(\alpha\) varies.
The central symmetric case and possible transitions in the saddle
structure merit separate analysis.

Another question is how far the proved critical strip can widen with
\(m\). The weighted localization proof identifies two competing
requirements: the endpoint amplitude must remain uniformly
expandable, and the tilted gamma tails must remain negligible.
An explicit maximal window, or several matched windows, would be
more informative than an unspecified two-variable asymptotic.

\subsection{A structural conjecture for all correction coefficients}
The construction in Theorem~\ref{moment:thm:all-orders} initially
produces rational functions of \(t\). Exact simplification removes
every negative power of \(t\) from \(C_1\) and \(C_2\). The following
is an exploratory coefficient identity, with evidence currently
limited to those two nontrivial orders.

\begin{conjecture}[Polynomial correction structure]
\label{future:conj:polynomial}
For every \(j\ge1\), the coefficient defined by
\eqref{moment:coefficient-algorithm} belongs to
\[
 \mathbb Q[\gamma,\gamma^{-1},\zeta(2),\ldots,\zeta(j+2)]
       [t,\lambda],
 \qquad
 \deg_t C_j\le j,\quad \deg_\lambda C_j\le2j,
 \qquad C_j(t,0)=0.
\]
\end{conjecture}

The theorem already proves the growth bound \(C_j=O(t^j)\) for
bounded \(\lambda\); the additional assertion is exact cancellation
of inverse powers and a uniform polynomial degree in \(\lambda\).
The diagonal residue identity
\eqref{moment:diagonal-generator} suggests a combinatorial mechanism.
Turning that suggestion into a coefficientwise argument would avoid
an invalid summation of the full divergent pole expansion.
A useful intermediate deliverable is an exact \(C_3\) or \(C_4\)
certificate obtained both from the gamma cumulants and from
finite residue diagonals.

\subsection{Lifting the depth congruences and resolving \texorpdfstring{\(S_6\)}{S6}}
The two congruences in Proposition~\ref{depthnew:prop:congruences}
can be lifted to complete functional identities by decomposing their
word residuals into shuffle products and transporting the endpoints.
An explicit lifted formula should contain every product and endpoint
constant, with branches fixed, and should be checked by differentiation
and one base value. This would make the formal reduction directly
usable for Gaussian special values.

At higher weights, compute the annihilators of the low-depth
\(\operatorname{GL}_2\)-orbit in the free Lie coalgebra.
The determinant-cube covariant at weight six is a particularly small
example. A representation-theoretic classification would replace
unstructured searches by a prediction of the missing directions,
while keeping evaluations at cyclotomic points separate.

For clarity, a concrete pre-existing target is retained below.
Let \(G=\beta(2)\), let
\(\beta(s)=\sum_{n\ge0}(-1)^n/(2n+1)^s\), and put
\[
 S_6=\sum_{n\ge0}\frac{(-1)^nH_n}{(2n+1)^6},
 \qquad g_{a,b}=\Im\Li_{a,b}(i,1).
\]
The pinned manuscript conjectures the weight-seven identity
\begin{align}
 S_6={}&\frac{722}{527}g_{6,1}
       +\frac{40}{527}g_{4,3}
       -\frac{128}{155}g_{2,5}
       +\frac{15191}{28569600}\pi^7\notag\\
      &-\frac3{124}G\zeta(5)
       -\frac{2373}{10540}\beta(4)\zeta(3)
       -2\beta(6)\log2.
 \label{future:eq:S6}
\end{align}
This article does not prove \eqref{future:eq:S6}.
The exact equivalence with the source's alternative coordinate
form is already established there. A substantive next step is to
find a convergent higher-weight analogue of the octahedral or
Cayley relation used for \(S_4\), or to reduce the residual in a
larger rigorously specified iterated-integral alphabet.
Repeating integer-relation searches in equivalent coordinate
systems supplies no missing analytic law.

\section{Reproducibility and integration}
\label{repro:sec}
The archive contains the full LaTeX source, the compiled article,
standalone figures, executable verification programs, their recorded
outputs, source provenance, and a proposed manuscript correction patch.
The four main mathematical sections can be integrated independently;
their labels have separate prefixes. The boundary section depends on
the signed-kernel section.

The golden certificate uses only integer and rational arithmetic,
checking six generating relations, two restricted-base identities,
eighteen primitive prime-ideal witnesses, and the finite restriction
ranks. Its independent audit checks every smaller residue exponent,
not just the proper-power tests in the original script. The infinite
support theorem additionally uses the cited primitive-divisor theorem.

The depth certificate reconstructs the Lie polynomial and all six
letter substitutions. It verifies the determinant-one minor, the
two quotient identities, and all 320 nonempty binary shuffle products
at weight six. These are exact finite calculations.

The moment program derives \(C_1,C_2\) by symbolic algebra and
checks the gamma moments about the mode. Its numerical part compares
direct \(x\)-quadrature with the transformed \(y=-\log x\) integral
in eight cases at 60 decimal digits of working precision.
The geometric diagnostics compare two kernel integrals, test the
critical density and deficit, and locate interior-disk zeros,
including the exact circular example. The boundary diagnostics
are documented alongside the additional local expansion.
All floating-point outputs are labeled as numerical checks, not
interval certificates or equality proofs.

The package's \path{README.md} gives the build and replay commands;
\path{integration/README.md} maps the sections to the pinned
manuscript. The exact scripts require only the Python standard
library; symbolic and numerical checks declare their additional
dependencies. Source hashes and the pinned commit make the
editorial patch reviewable against an unambiguous baseline.

The intended integration preserves the distinction between proved
theorems, formal quotient statements, established external inputs,
and the explicit conjectures in Section~\ref{future:sec}.
The latter should remain in the research agenda until an appropriate
proof or counterexample is supplied.
